\section{The normalized disk vacuum}\label{sec:vacuum}

The cylinder estimate used one differentiated row. For repeated rows with different boundary twists, we need normalized right and left vacuum caps. This section constructs the right cap as a simultaneous fixed vector of the spin rows, determines its exact polynomial denominator, and fixes the normalization under fusion. The left cap will then follow by transpose and spin reversal in Section~\ref{sec:site-moments}.

Polynomial exchange equations and special-rapidity reductions have earlier counterparts in the dense $O(1)$ sum rule of Di Francesco and Zinn-Justin~\cite{DiFrancescoZinnJustin2005} and in the dilute $O(1)$ ground-state work of Garbali and Nienhuis~\cite{GarbaliNienhuis2017Ground,GarbaliNienhuis2017Sum}, which gives a first normalization recurrence in a different convention. Our null-loop vector has different boundary data, and we prove its normalization directly below.

\subsection{The disk matching space}

We use the spin transfers \(T^\zeta(u;\mathbf b)\) of Section~\ref{sec:local}, with their labelled and possibly inhomogeneous site parameters, and take \(\zeta=i=p^{-2}\). The annular diagram space of Section~\ref{sec:cylinder} retains the homotopy of each arc around the cylinder. The space used here is instead a subspace of \(V^{\otimes N}\) defined by ordinary disk matchings, with no annular winding data.

For each noncrossing partial matching of the ordered sites in a disk, form a tensor by placing \(|0\rangle\) at each unmatched site and inserting \(C\) at the left and right endpoints of each arc, in that order. Let \(D\) be the span of these tensors. Every vector in \(D\) has total spin zero. The matching tensors are linearly independent: first separate them by their occupancy sets, then test a matching against the spin word that assigns \(+\) to each opening endpoint and \(-\) to each closing endpoint. Any matching compatible with that word has at least its word height many unclosed arcs at each cut, with simultaneous equality only for the chosen matching. Ordering by these heights gives a triangular test with nonzero diagonal. We call the coordinate on \(|00\ldots0\rangle\) the empty coefficient.

The adjacent operators \(\check R\) preserve \(D\). Their nonempty contractions use \(C^t\): this kills an arc \(C\), and otherwise joins its two partners with the prescribed phases, since cup-cap straightening uses \(p^m p^{-m}=1\). New arcs are inserted as \(C\). The last arc therefore cannot be capped off with nonzero weight; because the empty-to-empty entry is 1, \(\check R\) also preserves the empty coefficient. Twisted left rotation preserves matching tensors as well, by \(p^m\zeta^m=p^{-m}\).

The full row has the same two properties. To express its twisted trace using disk matchings, append an auxiliary and a spectator in \(|00\rangle+C\) at the right. Pass the auxiliary from last to first through the sites using adjacent \(\check R\)'s, and then contract the first and last exterior legs, namely the auxiliary and spectator, by \(\langle00|-C^t\). This gives exactly \(T^\zeta\), because \(-p^{2m}=\zeta^m\) for \(m=\pm1\). If the occupied exterior legs are joined to each other, their contraction is zero. If they are joined inward to partners, the contraction forms a new arc whose left spin is \(-m\), with phase \(-p^{3m}=p^{-m}\). Thus the output remains in \(D\). Again a last arc cannot disappear with nonzero weight, so only a completely empty input contributes to the empty output, with coefficient 1. Let \(D'\subset D\) be the invariant subspace of vectors with empty coefficient zero.

\subsection{Construction of the generic fixed vector}

We claim that for generic site parameters there is a unique \(\Psi_N(\mathbf b)\in D\) with empty coefficient 1 such that \(T^\zeta(u)\Psi_N=\Psi_N\) for every regular \(u\). Its components are rational in the half variables \(y_j=e^{ib_j}\). Set \(\Psi_0=1\) and \(z_j=y_j^2\). The proof finds finitely many row arguments for which a linear combination of \(T^\zeta-I\) is invertible on \(D'\); the empty normalization then gives the fixed vector by linear equations.

For that purpose, let the sites have very separated scales in \(\mathbb C^*\): successive ratios of their moduli, in increasing order, tend to infinity. Evaluate the row at intermediate scales \(z=e^{2iu}\), one in each gap, with its ratios to the scales on either side tending to their respective extremes. For \(z/z_j\to\infty\), the local weights satisfy \(a,b\to0,\ h\to1,\ d\to q^3,\ e\to q^{-3}\). The limit preserves occupancy on both legs, and the auxiliary spin can only increase or stay fixed: on \(|1,-1\rangle\), the opposite-spin off-diagonal weight is \(q^3+p^{-2}q^{-3}=0\). Its diagonal entries are \(q^{3s_a s_j}\). For \(z/z_j\to0\), the limits of \(d,e\) are exchanged; the auxiliary spin can only decrease or stay fixed, and the diagonal entries are \(q^{-3s_a s_j}\).

We will also need the same tests when the sites form nonempty blocks of comparable scales, with successive blocks very separated. Place the test scales between those blocks. The limiting rows preserve every site's occupancy and are simultaneously triangular on the full spin-zero sector, in any linear ordering compatible with the componentwise order of the block prefix sums. The blocks are ordered by increasing scale, which need not agree with their order around the disk. For a cut before the test gap, every spin in its prefix can only decrease. For a cut after the gap, every spin in the complementary suffix can only increase, so total-spin conservation again makes the prefix sum nonincreasing. If every block sum is preserved, all changes within a block have the same sign and hence every site spin is unchanged.

At the cut after block \(k\), let \(S_k\) be the spin in the smaller blocks. A contribution on the diagonal has constant auxiliary spin, so the diagonal entry of the limiting row is \(1+i q^{6S_k}-i q^{-6S_k}\). All the tested diagonal entries can equal 1 only when every block spin is even. Consequently, on every occupancy sector containing an occupied singleton block, a generic linear combination of the limiting \(T^\zeta-I\) is invertible, with one choice working for all such sectors. This is the interblock exclusion we shall use below.

For a fully separated list, every occupied sector has an occupied singleton block. Only the empty tensor remains, so the same combination is invertible on \(D'\), also before the limit when the scales are sufficiently separated. For \(N=1\), \(D'=0\) directly. Choose one regular configuration where invertibility holds and keep its test arguments and coefficients fixed. Rationality then makes invertibility generic in the site parameters. The combination maps \(D\) into \(D'\) and is invertible on \(D'\); its kernel in \(D\) is therefore one-dimensional and contains exactly one vector of empty coefficient 1. Since all rows commute with the combination and preserve the empty coefficient, this vector is fixed by every row. Solving the finite linear equations proves rationality in \(\mathbf y\) as well as the claimed uniqueness.

The row identities of Section~\ref{sec:local} now transfer to the normalized vector by uniqueness. It is invariant under a common translation of the site parameters, covariant under a period shift at one site by the sign of occupation there, and covariant under twisted cyclic rotation. Its adjacent exchange equation is
\[
 \Psi_N(\ldots,A,B,\ldots)=\check R_{j,j+1}(A-B)\Psi_N(\ldots,B,A,\ldots).
\]

\subsection{Counting the poles}

Existence and uniqueness determine the vacuum as a rational vector. To use its entries in the residue formula, we need its least common denominator, including the multiplicities of its factors. We first count poles while varying one site. The endpoint limits will then determine the leading coefficients, and symmetry and total degree will leave the fusion factors to be fixed.

For each component, divide by \(\prod_{j:\,\text{occupied}}y_j\). Sign covariance makes the result rational in \(\mathbf z\). Let \(\mathscr D_N(\mathbf z)\) be the least polynomial common denominator of these rational components on \((\mathbb C^*)^N\), chosen with no coordinate factor. It is unique up to a nonzero constant. We will prove that
\[
 \deg_{z_j}\mathscr D_N=N-1\quad(1\le j\le N),\qquad
 \deg\mathscr D_N=\frac{N(N-1)}2,
\]
that \(\mathscr D_N\) is symmetric and homogeneous, and that none of the divisors \(z_B/z_A=q^{\pm2},q^{\pm3},1\) is a factor. These divisor assertions concern generic points of each divisor; their intersections require separate justification when used.

Fix all sites except \(B\) and vary \(z_B\). Choose the fixed parameters generically, outside the loci where specialization creates cancellation, a degree drop, or a new coordinate factor, and then let their scales become very separated within this set. Such choices form a Zariski open set: a one-variable gcd identity keeps the reduced common fractions reduced after specialization. At a pole of the normalized vector, no combination of regular row transfers minus identity can be invertible on \(D'\), since the normalized linear equations would otherwise solve holomorphically there.

Consider any sequence of such poles as the fixed scales become more separated. After passing to a subsequence, the scales separate into singleton blocks except possibly for one pair \(A,B\) with comparable modulus. If no such pair remained, the fully separated tests would give invertibility on \(D'\) and rule out the poles. We may therefore assume that \(z_B/z_A\) has a nonzero finite limit, and pass to further subsequences to fix the square-root ratios. The interblock exclusion rules out all occupancy sectors with an occupied site outside this pair.

It remains to examine the pair. Evaluate additional rows at arguments for which \(e^{iu}/y_A\) is fixed and generic. In the limit these rows preserve the outside occupancies; on the subspace with empty spectators they reduce to the two-site transfer. The outside-occupancy decomposition also splits \(D\) and \(D'\). Hence degeneracy on \(D'\) requires a simultaneous eigenvalue 1 on the two-site \(D'\), which is the line spanned by \(C\). This follows either by commutation or by taking a generic combination of all the tests.

\subsection{Localization and existence of the simple poles}

Call the earlier site of the pair in disk order \(F\), and the other \(G\). In the two-site problem, evaluation at the first site \(u=b_F\) gives \(R_{FG}(b_F-b_G)\zeta^{s_F}\). On the disk matching space this acts by
\[
 C\mapsto d(b_F-b_G)C,\qquad |00\rangle\mapsto |00\rangle+b(b_F-b_G) C.
\]
At a regular first-site evaluation, subtracting the numerator and denominator of \(d(t)\) shows that \(d(t)=1\) requires \(\sin(2t)=0\). The case \(t=0\bmod\pi\) is excluded by evaluating instead at \(u=b_F+l\). Up to the sign gauges, the eigenvalue on \(C\) is then \(2x^4\ne1\): at \(l\) we have \(e=0,\ h=d=x^2\), and for outgoing pair spins \(m,-m\), the vacant traced auxiliary contributes \(p^m h^2\), while the occupied ones contribute \(\zeta^{-m}p^{-m}d^2=p^m d^2\).

This test also excludes a limiting pair for which the first-site evaluation itself has a denominator pole. Multiply the two-site operator by its two scalar sine denominators. At \(u=b_F\), its restriction to \(D'\) minus the cleared identity is nonzero at any such pole, because the numerator of \(d\) is nonzero there. It is consequently nonzero at some nearby regular row argument for that limiting pair. The only remaining possibility is \(t=\pi/2\bmod\pi\), or \(z_B/z_A\to-1\).

We now count the poles near each possible limit \(-1\). Take the interblock tests together with the first-site test \(u=b_F\), allowing that argument to move with the varying site. On a small disk in \(z_B/z_A\) about \(-1\), the transfers are analytic on each square-root sheet for sufficiently large separations and converge uniformly to the limiting operators. On \(D'\), those limits preserve the outside occupancies. Choose a generic linear combination of their differences from identity. On sectors with occupied spectators, the interblock exclusion makes it invertible at the center; the pair test can be given a sufficiently small nonzero coefficient to preserve this. On the remaining one-dimensional sector, the interblock differences vanish and the pair test has the simple zero of \(d-1\). The determinant of the combination on \(D'\) therefore has one simple zero at the center in the limit. The argument principle and the normalized linear equations show that, for sufficiently separated scales, the vector has at most one pole in the disk, and that pole is simple.

There is in fact a pole in each disk. On the whole spectator-empty pair sector in \(D\), the interblock transfers are the identity: both pair spins lie on the same side of a test, so neither can change, and their sum is zero. On a sufficiently small circle about \(-1\), the normalized vector therefore converges to the unique normalized two-site solution. Its coefficient of \(C\) is \(b/(1-d)\), which has nonzero residue. Contour integration forces a pole of the vector inside the circle. The two square-root signs give equivalent gauges. Passing to subsequences makes the argument apply along any sufficiently separated sequence, since the limiting tests do not depend on the extreme ratios.

There is thus exactly one simple pole near each antipodal value \(z_B=-z_A\), for the \(N-1\) fixed sites \(A\), and none away from those disks. Generic specialization gives \(\deg_{z_j}\mathscr D_N=N-1\) for every \(j\). The same tests show that \(\mathscr D_N\) has none of the divisors \(z_B/z_A=q^{\pm2},q^{\pm3},1\): make this pair comparable and the other sites separated, and use the pair tests away from \(-1\). They also show that a combination is invertible on \(D'\) at generic points of each such divisor. For \(N=1\), the degree and these assertions hold directly. The absence of the equal-site divisor is only a generic statement on that divisor; it does not by itself establish regularity when all sites coincide. Section~\ref{sec:physical-amplitude} proves that regularity for the sufficiently large sizes used below, from the projected annular row and its spectral gap.

\subsection{Endpoint limits, symmetry, and degree}

The pole count controls finite nonzero site values. We next let one site tend to zero or infinity while the others remain generic and fixed. In either limit, \(\Psi_N\) is bounded and converges to the embedded \(\Psi_{N-1}\) with that site vacant.

To prove this, evaluate one row at an intermediate scale separating the single departing site. Its triangular limit preserves occupancy and excludes an occupied singleton. Other fixed generic row arguments also preserve that site's occupancy in the limit and reduce on its vacant block to the shorter transfers. Together they give an invertible combination on \(D'\). On the vacant-site spin-zero block, the intermediate-scale test is the identity: monotonicity prevents any other site from changing, and the relevant prefix spins vanish. Thus the normalized solution of the limiting equations is precisely the embedded shorter vacuum. Invertibility gives boundedness and convergence to that solution. Throughout these arguments, the rows are evaluated at regular arguments before and in the indicated limits, and degeneracy is tested by finite linear equations.

Exchange and inversion now imply that \(\mathscr D_N\) is symmetric. Away from the exceptional ratio divisors, the exchange matrix is regular and invertible, so it preserves all polar divisors and their multiplicities. The exceptional divisors themselves are absent by the pole tests. Primitivity therefore makes a permutation act on \(\mathscr D_N\) by a scalar; an alternating sign is impossible because the equal-site divisor is absent. A common translation of the site parameters gives a common dilation of \(\mathbf z\). Translation invariance of the vector and primitivity similarly make \(\mathscr D_N\) homogeneous.

Its total degree is \(M_N=N(N-1)/2\). Choose rationally independent real slopes \(w_j\) in any increasing order and let \(z_j=\exp(Lw_j)\), \(L\to\infty\). Nonzero exceptional polynomials arising in fixed-site specializations stay nonzero for large \(L\). The root count and locations then give
\(z_j\partial_{z_j}\mathscr D_N/\mathscr D_N\) tending to the number of smaller sites. On the other hand, rational independence makes one monomial of \(\mathscr D_N\) uniquely dominant along this path, and these limits are its exponents. Their sum is \(0+1+\cdots+(N-1)=M_N\), which is the total degree by homogeneity.

Finally, regard \(\mathscr D_N\) as a polynomial in one site and pass to the endpoint limit in the divided components that are vacant at that site. Its leading coefficient times each limiting reduced component is a Laurent polynomial. Thus the leading coefficient clears every reduced denominator, and minimality of \(\mathscr D_{N-1}\), with no coordinate factors, makes it divisible as a polynomial by \(\mathscr D_{N-1}\). The total degree makes the quotient a constant. Normalize \(\mathscr D_0=\mathscr D_1=1\), and recursively choose this leading coefficient to be \(\mathscr D_{N-1}\). Symmetry makes the choice independent of the site.

\subsection{Fusion normalization}

With this normalization, define
\[
 Z_N=\mathscr D_{N}(\mathbf z)\prod_j y_j^{-(N-1)},\qquad \mathcal U_N=Z_N\Psi_N .
\]
Both are Laurent polynomials in \(\mathbf y\). For adjacent slots \(d_0+rl,d_0\), \(r=2,3\), let \({\rm red}\) denote the shorter list in which the pair is replaced by \(d_0+l\) for \(r=2\), or deleted for \(r=3\). We claim the recurrences
\begin{equation}\label{eq:vacuum-fusion}
 \Psi_N|_{\rm pair}=E_r\Psi_{\rm red},\quad
 Z_N|_{\rm pair}=2\cos(rl)\prod_{j\ {\rm other}}\prod_{h'=1}^{r-1} 2\cos(b_j-d_0-h'l)\ Z_{\rm red}.
\end{equation}
Here \(E_r\) is inserted in the indicated positions. The first formula follows at generic points of the pair divisor: \(E_r\) maps matching tensors into \(D\), preserves the empty normalization, and intertwines the transfers, so generic uniqueness applies there.

To prove the second formula, first locate the extra zeros on a pair divisor. All specializations of Laurent polynomials in this argument take place on the torus, with a chosen square-root sheet for the half variables. An ordered adjacent triple \(c+6l,c+3l,c\) forces \(\mathcal U_N=0\): polynomial continuation from the two pair divisors puts it in both \(E_3\) images, and these images have zero intersection because the two-site vector has matrix rank three. This wheel zero persists in cyclic orders with generic spectators inserted, by moving those spectators away through regular exchanges and then rotating. It also persists modulo \(\pi\) by the sign gauges.

On the specified adjacent pair divisor, the wheel zero therefore occurs at every spectator value \(b_j=d_0+h'l+4l\bmod\pi\). For \(r=3\), these values complete either end of the ordered wheel. For \(r=2\), the cyclic triple is \(d_0+2l,d_0,d_0+5l\), which modulo the period rotates to \(d_0+8l,d_0+5l,d_0+2l\). Since \(E_r\) is injective, the first fusion formula and polynomiality of \(\mathcal U_N\) show that \(\mathscr D_N|_{\rm pair}\Psi_{\rm red}\) is regular on the torus and vanishes on each of these extra divisors. Thus \(\mathscr D_N|_{\rm pair}\) is divisible as a polynomial by
\[
 \mathscr D_{\rm red}\prod_{j\ {\rm other}}\prod_{h'=1}^{r-1}(z_j+q^{h'}z_{d_0}),\qquad z_{d_0}=e^{2id_0}.
\]
This divisibility includes multiplicities even if one of the extra factors coincides with a divisor of \(\mathscr D_{\rm red}\). Indeed, at any divisor the minimum component order of \(\Psi_{\rm red}\) is minus the order of \(\mathscr D_{\rm red}\), by minimality of the denominator and the empty normalization. Adjoining a dummy variable or taking square roots on the torus does not change that assertion. Regularity supplies the order from \(\mathscr D_{\rm red}\), and the wheel zero supplies the additional order.

The quotient after division has total degree 1 and no spectator degree, so it is proportional to \(z_{d_0}\). After centering, this gives the displayed recurrence up to a constant. Compare the highest terms in each spectator in succession. The leading coefficient of \(y_j^{N-1}\) in \(Z_N\) is \(Z_{N-1}\prod_{k\ne j}y_k^{-1}\), and the removed pair product \(e^{i(2d_0+rl)}\) equals the product of the half variables contributed by the replacement site, when present, and the \(h'\)-factors. The comparison reduces the constant to the two-site case, where symmetry and normalization give \(Z_2=2\cos(b_1-b_2)\). This proves the second recurrence in \eqref{eq:vacuum-fusion}.

\subsection{An explicit component}

The fusion recurrences also determine the component with two consecutive blocks of opposite spins. This component will fix the normalization of the cap moments in Section~\ref{sec:site-moments}.

For \(N=2m\) and \(\epsilon=\pm1\), write \(\epsilon^m(-\epsilon)^m\) for the spin word consisting of \(m\) copies of \(\epsilon\) followed by \(m\) copies of \(-\epsilon\). Its component is
\begin{equation}\label{eq:two-block-component}
 (\mathcal U_N)_{\epsilon^m(-\epsilon)^m}
 =(p^\epsilon\,2\cos 3l)^m\!
 \prod_{\substack{i<j\\\text{same spin}}}
 H(b_i-b_j),\qquad H(t)=(2\sin(2l-t))(2\sin(3l-t)).
\end{equation}
To prove the formula, use the scalar exchange weight \(d\) within each block, interchanging generic intermediate sites of that block. It gives one factor \(H\) for each same-spin pair, so all the indicated factors divide the component on the half-variable torus. At either endpoint, an occupied component of \(\Psi_N\) vanishes; sign covariance gives at least one half-variable order of vanishing. After multiplication by \(Z_N\), the highest and lowest degrees of this component of \(\mathcal U_N\) therefore lie between \(-(N-2)\) and \(N-2\) in each variable. The quotient by the displayed product is constant. Fuse the middle pair at \(d_0+3l,d_0\). The two cosine factors for each spectator are precisely the removed \(H\), so \eqref{eq:vacuum-fusion} determines the constant recursively. Twisted rotation gives the other cyclic versions.
